% Options for packages loaded elsewhere
\PassOptionsToPackage{unicode}{hyperref}
\PassOptionsToPackage{hyphens}{url}
%
\documentclass[
  10pt,
  a4paper,
]{article}
\usepackage{amsmath,amssymb}
\usepackage{iftex}
\ifPDFTeX
  \usepackage[T1]{fontenc}
  \usepackage[utf8]{inputenc}
  \usepackage{textcomp} % provide euro and other symbols
\else % if luatex or xetex
  \usepackage{unicode-math} % this also loads fontspec
  \defaultfontfeatures{Scale=MatchLowercase}
  \defaultfontfeatures[\rmfamily]{Ligatures=TeX,Scale=1}
\fi
\usepackage{lmodern}
\ifPDFTeX\else
  % xetex/luatex font selection
  \setmainfont[]{DejaVu Serif}
  \setsansfont[]{DejaVu Sans}
\fi
% Use upquote if available, for straight quotes in verbatim environments
\IfFileExists{upquote.sty}{\usepackage{upquote}}{}
\IfFileExists{microtype.sty}{% use microtype if available
  \usepackage[]{microtype}
  \UseMicrotypeSet[protrusion]{basicmath} % disable protrusion for tt fonts
}{}
\makeatletter
\@ifundefined{KOMAClassName}{% if non-KOMA class
  \IfFileExists{parskip.sty}{%
    \usepackage{parskip}
  }{% else
    \setlength{\parindent}{0pt}
    \setlength{\parskip}{6pt plus 2pt minus 1pt}}
}{% if KOMA class
  \KOMAoptions{parskip=half}}
\makeatother
\usepackage{xcolor}
\usepackage[left=15mm,right=15mm,top=20mm,bottom=20mm]{geometry}
\usepackage{longtable,booktabs,array}
\usepackage{calc} % for calculating minipage widths
% Correct order of tables after \paragraph or \subparagraph
\usepackage{etoolbox}
\makeatletter
\patchcmd\longtable{\par}{\if@noskipsec\mbox{}\fi\par}{}{}
\makeatother
% Allow footnotes in longtable head/foot
\IfFileExists{footnotehyper.sty}{\usepackage{footnotehyper}}{\usepackage{footnote}}
\makesavenoteenv{longtable}
\usepackage{graphicx}
\makeatletter
\def\maxwidth{\ifdim\Gin@nat@width>\linewidth\linewidth\else\Gin@nat@width\fi}
\def\maxheight{\ifdim\Gin@nat@height>\textheight\textheight\else\Gin@nat@height\fi}
\makeatother
% Scale images if necessary, so that they will not overflow the page
% margins by default, and it is still possible to overwrite the defaults
% using explicit options in \includegraphics[width, height, ...]{}
\setkeys{Gin}{width=\maxwidth,height=\maxheight,keepaspectratio}
% Set default figure placement to htbp
\makeatletter
\def\fps@figure{htbp}
\makeatother
\setlength{\emergencystretch}{3em} % prevent overfull lines
\providecommand{\tightlist}{%
  \setlength{\itemsep}{0pt}\setlength{\parskip}{0pt}}
\setcounter{secnumdepth}{-\maxdimen} % remove section numbering
\usepackage{sectsty}
\allsectionsfont{\raggedright}
\usepackage{lineno}
\linenumbers
\usepackage{setspace}
\doublespacing
\usepackage{fancyhdr}
\pagestyle{fancy}\fancyhf{}
\fancyhead[L]{\small Tidal potential on a coastal–inland clock baseline}
\fancyfoot[C]{\thepage}\setlength{\headheight}{14pt}
\usepackage{float}\floatplacement{figure}{H}
\setlength{\emergencystretch}{2em}
\ifLuaTeX
  \usepackage{selnolig}  % disable illegal ligatures
\fi
\IfFileExists{bookmark.sty}{\usepackage{bookmark}}{\usepackage{hyperref}}
\IfFileExists{xurl.sty}{\usepackage{xurl}}{} % add URL line breaks if available
\urlstyle{same}
\hypersetup{
  hidelinks,
  pdfcreator={LaTeX via pandoc}}

\author{}
\date{}

\begin{document}

\hypertarget{tidal-potential-measured-on-a-coastalinland-optical-clock-baseline}{%
\section{Tidal potential measured on a coastal--inland optical-clock
baseline}\label{tidal-potential-measured-on-a-coastalinland-optical-clock-baseline}}

{[}Author list, order and affiliation assignments to be confirmed by the
authors{]}

\textsuperscript{1} Hefei National Laboratory, University of Science and
Technology of China, Hefei, China

\textsuperscript{2} Shanghai Research Center for Quantum Sciences and
CAS Center for Excellence in Quantum Information and Quantum Physics,
University of Science and Technology of China, Shanghai, China

\textsuperscript{3} Hefei National Research Center for Physical Sciences
at the Microscale and School of Physical Sciences, University of Science
and Technology of China, Hefei, China

\textsuperscript{4} Innovation Academy for Precision Measurement Science
and Technology, Chinese Academy of Sciences, Wuhan, China

\textsuperscript{5} Jinan Institute of Quantum Technology and CAS Center
for Excellence in Quantum Information and Quantum Physics, University of
Science and Technology of China, Jinan, China

\textsuperscript{6} Time and Frequency Geodesy Center (TFGC), School of
Geodesy and Geomatics, Wuhan University, Wuhan, China

\hypertarget{abstract}{%
\subsection{Abstract}\label{abstract}}

Optical clocks measure changes in Earth's gravitational potential
through their relative rates.\textsuperscript{1,2} Resolving
neighbouring tidal components is difficult when observations are
interrupted and their frequency zero points are not transferred. Here we
resolve semidiurnal responses with principal frequencies only 0.78 µHz
apart on a coastal--inland clock baseline. Shanghai strontium and Wuhan
ytterbium clocks accumulated 545.9 hours over 353 days. The changing
lunar--solar phase provides an external reference across 54 segments
with independent frequency offsets. The measured spectrum has local
peaks near the principal lunar and solar frequencies. Joint estimation
of astronomical templates gives responses of \(1.00\pm0.13\) for M₂+N₂
and \(0.93\pm0.36\) for S₂+K₂ relative to the combined solid-Earth and
ocean-loading prediction (one statistical standard error). These results
demonstrate spectrally resolved measurements of a changing gravitational
potential with an optical-clock link operated intermittently.

\hypertarget{main-text}{%
\subsection{Main text}\label{main-text}}

Optical clocks measure gravitational potential through their relative
rates.\textsuperscript{1,2} Laboratory comparisons have resolved
height-dependent redshifts, and remote comparisons have supported
geodetic measurements between distant sites.\textsuperscript{3,4,5}
Stabilized optical-fibre networks extend these measurements across
regions without transporting the clocks.\textsuperscript{6,7}
Coordinated international comparisons now connect optical standards
across national borders, supporting their use in shared time
scales.\textsuperscript{8,9} The same links can repeatedly sample
potential differences between distant sites, complementing levelling
surveys and satellite gravity measurements.\textsuperscript{10,11}
Measuring a changing potential requires its temporal components to be
distinguished from fluctuations in the clocks and transfer system.
Astronomical tides provide a physical test of this capability because
their frequencies and phases are predictable independently of the clock
observations.\textsuperscript{12}

The principal lunar and solar semidiurnal tides are a demanding pair of
nearby signals. M₂ repeats every 12.42 h and S₂ every 12 h, giving a
frequency separation of only 0.78 µHz. Over a short comparison their
variations can look much alike. Their relative phase changes with date,
allowing observations distributed over different lunar--solar phases to
separate their responses. That astronomical phase remains defined
through interruptions, even when the frequency offset between observing
runs is unknown.

Here we use this external phase reference in a comparison between
coastal Shanghai and inland Wuhan. Across a 353-day span, 54 valid
segments provide 545.9 h of observations. Independent constants
represent their reference levels; the tidal phase is retained at every
actual observing time. We jointly estimate complete physical templates
for the lunar semidiurnal group M₂+N₂, the neighbouring group S₂+K₂ and
the approximately daily group K₁+O₁+P₁+Q₁. Internal amplitude ratios and
phases remain fixed by the astronomical model, so the fit propagates
overlap between the groups. K₂ contains lunar and solar forcing. The
spectrum displays local peaks near M₂ and S₂, and both fitted
semidiurnal responses agree with the combined solid-Earth and
ocean-loading prediction. The experiment thus tests spectral separation
in an intermittently operated clock link, on a baseline where loading is
a sizeable part of the signal.

\hypertarget{a-coastalinland-prediction}{%
\subsubsection{A coastal--inland
prediction}\label{a-coastalinland-prediction}}

The strontium clock in Shanghai and ytterbium clock in Wuhan are 692 km
apart, joined through Hefei by a 1,350-km phase-stabilized fibre link
(Fig. 1a).\textsuperscript{13} The summer observations overlap the
absolute frequency-ratio measurement reported by Hou and colleagues; the
present observable is the time-dependent tidal response. The two clock
systems have published systematic evaluations, which describe their
apparatus rather than the uncertainty of the interrupted tidal
estimate.\textsuperscript{14,15} We define a positive frequency-ratio
variation as Wuhan Yb running faster relative to Shanghai Sr.~In the
geodetic potential convention used here, the measured fractional
variation is
\(y=-[\delta W_{\mathrm{Wuhan}}-\delta W_{\mathrm{Shanghai}}]/c^2\).

Tides change the direct gravitational potential and deform the Earth,
moving the clocks and changing the Earth's own potential. Ocean tides
additionally redistribute water. Their contribution includes its
attraction, the Earth's deformation and the accompanying displacement of
the clock. We combine an IERS-based solid-tide prediction with FES2022
ocean elevations and a loading kernel containing all three loading
contributions (Methods).\textsuperscript{16,17,18,19} Shanghai lies by
the East China Sea, whereas Wuhan is inland. This geometry makes loading
a substantial part of the differential prediction.

The predicted M₂ amplitudes are \(3.43\times10^{-18}\) for the solid
contribution and \(1.42\times10^{-18}\) for loading. Their phases differ
by 19.8°, giving a combined amplitude of \(4.78\times10^{-18}\) (Fig.
1b). The loading amplitude is about 30\% of the combined M₂ amplitude;
their vector sum accounts for the phase difference. On the sampled
record, after removing segment means, loading has 33.1\% of the combined
prediction's weighted root-mean-square variation. These amplitude ratios
describe the model contribution, making this coastal--inland geometry a
useful setting for testing the combined potential. A separately
processed analysis allowing loading its own coefficient reported
\(0.89\pm0.55\), which does not resolve that coefficient (Methods). The
observable tested here is the combined potential.

The record spans September 2025 to August 2026, with 54 valid segments
in eight observing periods (Fig. 1c). Its 3,299 ten-minute means contain
545.9 h of accepted nominal one-second measurements. Observations occur
in eight calendar months; the annual span is not continuous operation.
Actual UTC times and all gaps are retained. The June--August
observations supply 273.9 h and sample a broad range of relative
lunar--solar phase, while the earlier periods add different phases and
operating conditions (Fig. 1d). Both the duration and placement of the
observations determine what can be separated.

\hypertarget{jointly-resolved-semidiurnal-responses}{%
\subsubsection{Jointly resolved semidiurnal
responses}\label{jointly-resolved-semidiurnal-responses}}

Fitting the three physical groups together gives \(A_L=1.00\pm0.13\) for
M₂+N₂ and \(A_S=0.93\pm0.36\) for S₂+K₂. One denotes the specified
combined prediction. Each of the 54 segments has its own constant
reference level, and the same noise weights are used throughout the
amplitude analysis. The coefficients are estimated from the original
ten-minute means. No empirical daily sinusoid is subtracted first, and
no model values replace missing measurements.

Figure 2 orders the entire record by M₂ phase. The observations, after
subtracting the fixed minor tidal remainder, and the unit lunar
prediction undergo the same projection against segment constants and the
other two fitted groups. Twelve phase bins then reveal the recurring
modulation. All individual points and the number supporting each bin are
retained. The propagated error bars account for the jointly estimated
nuisance terms; the twelve displayed means are a view of the full fit,
rather than twelve new independent observations.

Allowing a lunar component one quarter-cycle from the predicted phase
gives in-phase and quadrature responses of \(1.00\pm0.13\) and
\(-0.10\pm0.20\) (Fig. 3b). The phase is compatible with the
astronomical prediction. The jointly fitted daily group has
\(A_D=0.58\pm0.54\) (Fig. 3a). Its lower precision is already accounted
for through the joint design and covariance; it does not negate the
semidiurnal responses. It also does not provide a bound on an unmeasured
laboratory daily disturbance.

The reported errors are the larger of residual-scaled least-squares and
UTC-day-deletion standard errors. Their components, period-deletion
checks and conditional interval coverage are given in Methods. Joint
estimation preserves uncertainty from overlap among the physical groups.
Improving the daily response calls for longer uninterrupted clock
operation and sustained link comparisons with all valid observations
retained under data-quality rules fixed independently of the tidal fit.
An existing single-bin influence check is described in Methods.

\hypertarget{spectral-separation-through-the-observation-window}{%
\subsubsection{Spectral separation through the observation
window}\label{spectral-separation-through-the-observation-window}}

The spectrum presents the same observations in a complementary,
physically intuitive form (Fig. 3c). An equal-weight sinusoidal scan
with 54 segment constants has local maxima at 22.361 and 23.147 µHz,
close to the predicted M₂ and S₂ frequencies of 22.364 and 23.148 µHz.
The prescribed principal frequencies are 0.78 µHz apart, compared with
an inverse calendar span of 0.0328 µHz. Resolving this pair is the
demonstrated spectral capability; the interrupted observation window
determines the accompanying sidelobes. The left axis expresses spectral
density in gravitational-potential units and the right in
fractional-frequency units. A fractional variation of \(10^{-18}\)
corresponds to about \(0.09\,\mathrm{m^2/s^2}\) of potential difference.

Interruptions shape this spectrum. Sampling a narrow tidal line at the
observed times spreads part of its power into sidelobes, including the
neighbourhood of the other principal semidiurnal frequency.
Single-constituent injections through the actual sampling and scan give
8.03\% M₂-to-S₂ cross-power and 8.15\% in the reverse direction. These
numbers describe the sampling response, not random noise added by the
gaps. Conditional scans retain a local feature near each frequency after
projecting out the sine and cosine at the other (Extended Data Fig. 2).

The two features and the conditional scans support distinct
contributions in the principal lunar and solar semidiurnal
neighbourhoods. This attribution does not separately measure N₂ or K₂.
The response coefficients instead come from the simultaneous
astronomical groups, with their prescribed internal relationships and
noise weights. Subtracting the fixed unit physical prediction reduces
the spectral structure, as shown by the second curve. That curve is a
residual display of the same record, not another significance test.
Phase coverage, spectral attribution and grouped estimation answer
related but distinct questions.

\hypertarget{across-interruptions-and-through-the-fibre}{%
\subsubsection{Across interruptions and through the
fibre}\label{across-interruptions-and-through-the-fibre}}

The response is supported by observations on both sides of summer. With
the same segment model, the lunar coefficient is \(1.21\pm0.23\) before
summer and \(0.77\pm0.18\) in summer. Deleting each observing period in
turn gives point estimates from 0.885 to 1.093 (Fig. 4a). These
overlapping fits show the influence of each period; the headline result
uses the full record.

The choice of reference levels follows the experimental record. The
available interval and correction tables do not establish a calibrated
zero transfer across interruptions. We therefore give each valid segment
an independent constant. The model with eight period constants
reproduces \(0.987\pm0.123\) on exactly the same observations, templates
and weights. Its lunar result changes little, but its neighbouring-group
coefficient becomes \(1.129\pm0.288\). Residual structure and deletion
results for both designs are retained in Methods and Extended Data Fig.
1. Agreement with one is not used to choose between them.

A simultaneous fibre round-trip record provides a direct instrument
control on 246.6 h of summer observations, from 29 June to 25 August
2026. Clock and monitor means are formed from matching valid members.
Adding the monitor to the otherwise identical segment model changes the
lunar coefficient by \(0.0061\pm0.0063\) and the neighbouring
coefficient by \(0.0059\pm0.0145\) (Fig. 4b). This tests the monitored
transmission route on that overlap. Other local transfer branches and
clock effects require their own observations.

Daily temperature changes and operating schedules motivate environmental
checks. The separately filtered analysis reports a lunar coefficient of
\(0.98\pm0.13\) with two-site temperature regressors; its distinct
sampling and supporting measurements are identified in Methods. On the
present observations, adding one shared 24-h sine--cosine pair gives
\(1.04\pm0.13\). Arbitrary period-specific 12-h terms are retained as
sensitivity tests because, over short runs, they compete directly with
astronomical semidiurnal variation. They are not assigned an observed
laboratory amplitude.

The comparison resolves nearby gravitational signals from observations
separated by interruptions in clock and link operation. A phase
reference supplied by celestial motion connects the physical waveforms
across the observing dates, while segment constants accommodate
frequency levels that have not been transferred between runs. The
measured peaks and jointly estimated groups resolve neighbouring
semidiurnal responses without forcing the clock record onto a common
level. This allows an intermittently operated clock link to measure
distinct temporal components of the gravitational potential.

Our observations and contemporaneous European clock-network measurements
by Stahl and colleagues are independent, parallel
demonstrations.\textsuperscript{20} The European network tracks
composite tidal variations on different baselines; the present
coastal--inland comparison emphasizes separation of neighbouring
semidiurnal responses across interrupted records. Together they show the
scope of clock networks as instruments for time-dependent gravitational
potential. Longer coastal records and calibrated local transfer paths
could allow the solid and loading responses to be tested separately.

Clock networks connect this geophysical capability to the realization of
time itself. International comparisons already link optical standards
across countries, and intermittently operated optical clocks have
steered continuous time scales.\textsuperscript{8,21} A common time
scale requires clock frequencies to be referred to a consistent
gravitational potential; Earth monitoring treats changes in that
potential as the signal.\textsuperscript{9,10} Once a frequency-transfer
link is established, it can repeatedly compare distant endpoints during
operation without resurveying the intervening terrain. Such observations
complement levelling surveys and the standard monthly gravity products
from satellites.\textsuperscript{10,11} The temporal resolution
attainable for a particular signal depends on clock stability and
simultaneous operation. With improved continuity, calibrated reference
levels and colocated height measurements, future networks could extend
these differential observations to atmospheric and hydrological mass
changes.\textsuperscript{11,22} Resolving known tidal components across
interruptions tests the signal discrimination that those applications
require.

\clearpage

\singlespacing

\hypertarget{figure-legends}{%
\subsection{Figure legends}\label{figure-legends}}

\includegraphics{figures/figure1.pdf}

\textbf{Fig. 1 \textbar{} Geometry and changing lunar--solar phase.}
\textbf{a,} Shanghai Sr and Wuhan Yb, 692 km apart, with the 1,350-km
fibre connection through Hefei drawn schematically. \textbf{b,}
Predicted differential M₂ solid and ocean-loading phasors, with the
solid component defining phase zero; amplitudes are fractional clock
rates in units of \(10^{-18}\). \textbf{c,} Actual UTC sampling in eight
observing periods, with accepted nominal member-hours. \textbf{d,} The
S₂−M₂ astronomical phase at those same samples, including nodal
corrections. Summer periods P6--P8 are orange. Geography and phasors
describe the prediction; the lower panels describe the observation
window.

\clearpage

\includegraphics{figures/figure2.pdf}

\textbf{Fig. 2 \textbar{} The full record ordered by lunar phase.}
\textbf{a,} Twelve 30° M₂ phase means of all 3,299 observations, after
subtracting the fixed minor remainder and projecting against 54 segment
constants and the neighbouring semidiurnal and daily groups. Identical
weighted projection and averaging are applied to the unit lunar
prediction; the fitted curve uses the unbinned primary coefficient.
Error bars are one working standard error with the nuisance projection
propagated; phase means are correlated. \textbf{b,} Every projected
observation, with the full observed vertical range. \textbf{c,} Number
of rows in each bin. Data, weights and templates are the primary segment
design throughout. This is a display of the main observations, not an
independent experiment.

\clearpage

\includegraphics{figures/figure3.pdf}

\textbf{Fig. 3 \textbar{} Semidiurnal responses, phase and measured
spectral structure.} \textbf{a,} Joint group amplitudes with one
standard error, defined as the larger of scaled least-squares and
UTC-day jackknife errors. \textbf{b,} Lunar in-phase and quadrature
coefficients; the cross marks the unit prediction. Contours have
covariance radii one and two from scaled least squares, not separately
calibrated joint confidence levels. \textbf{c,} Equal-weight spectral
calculation around the principal M₂ and S₂ frequencies, separated by
0.78 µHz: 3,299 real UTC means, 54 segment constants and equal weights.
Blue shows the observed scan; orange follows subtraction of the fixed
unit total prediction. No gap is filled or compressed. Both axes show
the same nominal one-sided projection density per µHz, related by
\(S_W=c^4S_y\). Sampling sidelobes contribute about 8\% bidirectional
M₂/S₂ cross-power. Peak locations support the principal-frequency
attribution; group coefficients in \textbf{a} use the physical templates
and Allan-based weights. All panels use the same observation rows.

\clearpage

\includegraphics{figures/figure4.pdf}

\textbf{Fig. 4 \textbar{} Temporal support and a simultaneous
transmission control.} \textbf{a,} Lunar estimates from the full record,
the pre-summer and summer subsets, and deletion of each observing
period. Each fit uses segment constants and the primary noise weights,
with errors recalculated on its retained rows. The shaded band is the
full-record one-standard-error range. These fits share observations.
\textbf{b,} Changes caused by including the simultaneous round-trip
monitor, with paired UTC-day-jackknife standard errors. Both fits use
identical clock/monitor members: 1,482 means representing 246.6 h in
P6--P8, between 29 June and 25 August 2026. This control covers the
monitored transmission channel on its measured overlap; the principal
tide fit retains the entire 545.9-h clock record.

\clearpage

\doublespacing

\hypertarget{methods}{%
\subsection{Methods}\label{methods}}

\hypertarget{observations-frequency-sign-and-averaging}{%
\subsubsection{Observations, frequency sign and
averaging}\label{observations-frequency-sign-and-averaging}}

The Shanghai and Wuhan coordinates used by the frozen prediction are
31.130° N, 121.570° E and 30.529° N, 114.357° E. The measured ratio is
Wuhan Yb/Shanghai Sr.~With positive geodetic potential \(W\), its tidal
fractional variation is
\(-[\delta W_{\mathrm{Wuhan}}-\delta W_{\mathrm{Shanghai}}]/c^2\). The
period-relative beat conversion is
\(-d_p(f_{\mathrm{beat}}-f_{{\mathrm{ref}},p})/F_{1550}\), where
\(F_{1550}=1.933992\times10^{14}\) Hz and \(d_p=2\) in 2025 and 1 in
2026. Thus a 2026 beat deviation of 1 Hz has fractional magnitude
\(5.17\times10^{-15}\). The 2025 division factor is applied before
forming the small clock-rate variation.

The supplied valid-interval table defines 54 segments in 29 operating
groups and eight periods. Counter members are read across midnight
before binning. Beijing-time labels are converted to UTC. Summer group
corrections follow the decimal-arithmetic frequency-chain calculation;
early absolute correction entries remain incomplete. Their constant part
is represented by segment intercepts, while any time variation within a
segment remains a possible systematic contribution. Fitted intercepts
are never applied to join the runs for long-term stability estimation.

The membership audit contains 1,966,908 nominal one-second rows before
the minimum-bin-support cut. Three boundary rows reconcile its
difference from the printed interval total. Means are formed in nominal
600-s bins separately within each segment and require at least 180
retained members. The 3,299 accepted means contain 1,965,327 members, or
545.924 h. There are 3,291 unique bin centres: eight centres contain
separate partial means from different segments, with disjoint members.
The 353.39-day span contains observations on 43 UTC dates. The inherited
valid mask is retained; this work neither applies an additional
triangular filter nor selects points using tidal residuals. Exact
raw-file memberships and hashes are retained in the controlled archive.

\hypertarget{what-the-constituent-names-mean}{%
\subsubsection{What the constituent names
mean}\label{what-the-constituent-names-mean}}

Semidiurnal tides produce approximately two cycles per day. M₂ is the
principal lunar constituent; N₂ describes an additional lunar variation
associated with orbital ellipticity. S₂ is the principal solar
semidiurnal constituent, whereas K₂ contains lunar and solar forcing.
The names identify astronomical frequencies, not separate instruments or
independent experiments. The diurnal group contains K₁ (lunisolar), O₁
(principal lunar), P₁ (principal solar) and Q₁ (lunar elliptic). Their
periods are listed in Extended Data Table 3.\textsuperscript{23}
Internal relative amplitudes, phases and nodal modulation are fixed
within each physical group. A group response therefore tests that
complete specified pattern.

\hypertarget{combined-solid-and-loading-potential}{%
\subsubsection{Combined solid and loading
potential}\label{combined-solid-and-loading-potential}}

The physical target includes the direct astronomical potential, the
Earth's induced potential and the effect of clock displacement. The
selected frozen solid sequence implements these through the IERS radial
displacement \(u\) and a fixed effective-potential conversion,

\[y_{\mathrm{solid}}=-\frac{1+k_2-h_2}{h_2}\frac{g(u_{\mathrm{Wuhan}}-u_{\mathrm{Shanghai}})}{c^2},\]

with \(h_2=0.6078\), \(k_2=0.3010\) and \(g=9.794\,\mathrm{m/s^2}\).
This is a displacement-based approximation to the combined solid
potential, not the alternative frequency-resolved direct-potential
implementation. The same conversion also acts on degree-three and
frequency-dependent displacement contributions in the supplied sequence.
A previous implementation comparison found an approximately 18\% K₁
amplitude deficit relative to a potential-weighted reference. The daily
response is accordingly not used to infer a Love number. Uncertainty in
the solid model remains separate from the quoted statistical errors.

For loading, FES2022 sea-surface elevations \(\zeta\) are convolved with
the effective-potential kernel

\[\delta W_{\mathrm{load}}=\int \rho_w\zeta(\mathbf{r}',t)G_W(\psi)\,\mathrm{d}A,\qquad
G_W(\psi)=\frac{ga}{M_E}\sum_n(1+k'_n-h'_n)P_n(\cos\psi).\]

The three terms represent direct water attraction, induced potential and
clock displacement. PREM load Love numbers are retained through degree
696 with the archived high-degree continuation. Eighteen loading
constituents are synthesized with their astronomical nodal corrections.
The complex convention is \(f\,\mathrm{Re}[C\exp(-i\phi)]\); reversing
the exponential sign requires conjugating \(C\). The present analysis
reads the frozen constituent and time-series tables rather than
substituting an older loading sequence.

The eight named constituents are grouped as described above. Minor terms
remain once as a fixed remainder \(r_0\). On the actual weighted
segment-centred sampling, solid, loading and total root-mean-square
variations are respectively 2.928, 1.328 and \(4.011\times10^{-18}\).
Their ratio 0.331 is an amplitude-norm ratio, not an additive fraction
of total power. The separately reported \(0.89\pm0.55\) loading
coefficient belongs to the alternative directly modelled, filtered
reduction described below; it has not been recalculated as a separate
coefficient of the present estimator.

\hypertarget{joint-estimator-and-zero-point-choice}{%
\subsubsection{Joint estimator and zero-point
choice}\label{joint-estimator-and-zero-point-choice}}

We fit

\[y_i=c_{s(i)}+A_Ls_{L,i}+A_Ss_{S,i}+A_Ds_{D,i}+r_{0,i}+\epsilon_i,\]

with 54 segment intercepts and three physical response coefficients.
Weighted least squares uses the frozen Allan-based uncertainties
\(\sigma_i\), with \(W_{ii}=1/\sigma_i^2\). A period noise coefficient
\(\alpha_p\) is estimated from within-segment overlapping Allan
deviations at 256, 512, 1,024, 2,048 and 4,096 s, after removing the
fixed unit tide. No Allan window crosses a segment boundary. The bin
uncertainty is \(\sigma_i=\alpha_{p(i)}/\sqrt{n_i}\), where \(n_i\) is
its retained member count. This is a working noise model, not stability
calculated after joining fit-adjusted levels. Independent QR replay
reproduces the eight-period design on the same rows before changing only
its intercept grouping. All comparisons therefore hold sampling,
correction values, physical templates and weights fixed.

The inspected interval and correction records contain no documented
calibrated zero transfer for the 46 adjacent segment boundaries that
share an eight-period intercept. They give valid times, but not a signed
boundary transfer with uncertainty or a lock/relock history establishing
unchanged reference levels. This motivates the segment design. The
eight-period model's weighted residual sum of squares is 3,781.58; the
segment model's is 3,713.64, with 46 additional intercept parameters.
Their unweighted residual root-mean-square values are 24.114 and
\(23.934\times10^{-18}\). Reduced residual sums of squares are 1.1501
and 1.1455. The numerical improvement alone does not establish a
physical jump at each boundary.

With period constants, weighted segment-mean residuals range from
\(-6.75\) to \(12.00\times10^{-18}\). Adjacent standardized residual
correlations, restricted to the same segment and separations of at most
900 s, are −0.127 and −0.148 for the period and segment designs.
Deleting a period gives lunar point-estimate ranges 0.877--1.052 and
0.885--1.093, respectively (Extended Data Fig. 1). These checks retain
substantial lunar support under either reference-level assumption. They
do not supply missing calibration logs. Coefficient proximity to one is
not a model-selection rule.

\hypertarget{statistical-errors-and-conditional-coverage}{%
\subsubsection{Statistical errors and conditional
coverage}\label{statistical-errors-and-conditional-coverage}}

For each group the reported standard error is the larger of
residual-scaled least squares and leave-one-UTC-day-out jackknife
errors. Deletions remove a calendar day, refit all remaining
coefficients and drop empty intercept columns. The jackknife variance is
\((G-1)G^{-1}\sum_g(\hat A_{(-g)}-\overline A)^2\). Period-deletion
errors are recorded separately. Extended Data Table 1 gives all
components and the same-data intercept comparisons. Observations are
repeated measurements of one clock pair, not thousands of independent
experiments.

New conditional simulations apply both designs to the same 3,299-row
sampling and heteroscedastic uncertainties. Each of four scenarios has
1,000 realizations: independent Gaussian noise; equal variance
contributions from Gaussian white noise and a 6-h Ornstein--Uhlenbeck
process within segments; the corresponding 24-h process within periods;
and standardized Student-\(t_5\) noise. The overall scale is the
observed segment-model residual scale; all true physical response
coefficients are one. Each realization recomputes both components of the
standard-error rule. Seed 2026110617 fixes the simulation. Coverage of
\(\hat A\pm1.96u\) is 95.1--97.4\% for the segment-model lunar
coefficient and 95.4--96.8\% for its neighbouring coefficient; Monte
Carlo standard errors are about 0.4--0.7 percentage points (Extended
Data Table 2). These are tests under declared noise laws, not a
calibration against every possible disturbance.

The phase-relaxed model adds the lunar quadrature while retaining the
prescribed M₂/N₂ relationship. Covariance contours use scaled least
squares. Subset and deletion fits refit the same model on their retained
rows. No shared-data estimates are averaged, and no half-daily
coefficient is converted into a daily upper limit.

\hypertarget{phase-means-spectral-density-and-sampling-sidelobes}{%
\subsubsection{Phase means, spectral density and sampling
sidelobes}\label{phase-means-spectral-density-and-sampling-sidelobes}}

For the phase plot, \(N\) contains segment constants and the two
non-lunar groups. We project both \(y-r_0\) and the unit lunar
prediction with \(M_N=I-N(N^TWN)^{-1}N^TW\), and average in twelve fixed
M₂ phase bins. If \(B\) is weighted bin averaging, the mean covariance
is \(BM_NC_yM_N^TB^T\), where \(C_y\) is the scaled working observation
covariance. It retains correlations between means and propagates
nuisance estimation. The full unbinned response remains primary.

Figure 3c preserves the independently specified equal-weight,
54-intercept scan of the unchanged observed means, rather than replacing
it with the amplitude fit's weights. Every trial sine--cosine pair and
the data are projected against the same constants. For a reduction in
squared residuals \(\Delta\mathrm{RSS}(f)\) with dimensionless \(y\),
the nominal one-sided projection density is
\(S_y(f)=600\,\mathrm{s}\,\Delta\mathrm{RSS}(f)\) per Hz, or \(10^{-6}\)
times that value per µHz. Potential density is \(S_W=c^4S_y\); amplitude
density is \(c^2\sqrt{S_y}\). Uniform-grid checks reproduce the
conventional one-sided Fourier density. For interrupted data this
normalization is a stated projection density, not a Parseval-calibrated
variance over a filled year.\textsuperscript{24}

The frequency step is 0.002729 µHz; the inverse calendar span is
0.032752 µHz. Fine plotting steps do not improve physical resolution.
M₂-only and S₂-only astronomical injections use the same actual times,
nodal factors and equal-weight segment projection. Power at the opposite
principal frequency, divided by power at the injected frequency, gives
0.08027 and 0.08149. Conditional scans project out a freely phased
sinusoid at the other frequency before scanning the target
neighbourhood. Amplitude interference means that such corrections act on
sampled waveforms or complex coefficients, not by subtracting scalar
power spectra. Gaps alone do not predict the realization of measurement
noise. Residual density can describe remaining measurement fluctuations,
but includes possible model mismatch and cannot be called a pure
instrumental noise floor.

\hypertarget{instruments-environmental-checks-and-the-alternative-reduction}{%
\subsubsection{Instruments, environmental checks and the alternative
reduction}\label{instruments-environmental-checks-and-the-alternative-reduction}}

An existing fixed-weight single-bin diagnostic in the eight-period
reduction examined the pre-identified extreme mean at 2026-07-07 01:25
Beijing time (segment 46). Omitting that mean changed the daily
coefficient from 0.678 to 0.721, a shift of 0.043 compared with its
0.359 standard error. The lunar coefficient changed by 0.039. This
archived diagnostic kept the template and weights fixed; it did not
recalibrate uncertainty. The point remains in the present 3,299-row
sample. Its influence does not explain the weaker daily constraint or
justify selecting observations to improve agreement with the tide.

The matched monitor calculation uses 1,482 rows with 887,882 nominal
members, or 246.634 h, in P6--P8. Both baseline and augmented fits have
identical members, segment labels and weights. Paired day-deletion
replicates propagate uncertainty in their difference. The lunar,
neighbouring and daily coefficient changes are respectively
\(0.00606\pm0.00627\), \(0.00594\pm0.01453\) and \(0.00791\pm0.01388\).
The archive does not provide that paired control for the pre-summer
periods or every laboratory transfer branch.

The alternative v68 reduction reports \(A_L=0.97\pm0.14\) from
approximately 544 h in 54 segments, with a 1,200-s triangular filter
sampled every 600 s, equal weights, a direct-potential solid model and a
continuous-run circular-residual-shift error. It reports \(0.98\pm0.13\)
with two-site ERA5 temperature regressors and a 474-h round-trip overlap
changing lunar and whole-template coefficients by less than 0.005. These
are author-reported manuscript results, not the present replay. The
stated 98-h Sr chamber record is a distinct thermometer measurement;
ERA5 2-m air temperatures are environmental proxies. Exact membership
for the 474-h control and temperature-regression support must be
supplied for an independent replay. These checks are not combined with
the current uncertainty.

On the selected data and model, a common 24-h sine--cosine pair gives
\(A_L=1.036\pm0.133\); separate 24-h pairs in each period give
\(1.003\pm0.136\). These are mechanism-motivated daily sensitivity
designs, not measurements of an independently observed thermal
amplitude. Free 12-h pairs in each period give \(0.751\pm0.242\) and a
very poorly constrained neighbouring coefficient; the design competes
with astronomical semidiurnal variation within short periods. It is
retained as a sensitivity boundary, with no assumed laboratory
disturbance assigned to it. A sequential residual diagnostic must
project its test basis through the same physical-fit operator and
propagate uncertainty.

\singlespacing

\hypertarget{data-availability}{%
\subsection{Data availability}\label{data-availability}}

The author review package contains the 3,299 derived means, fixed
templates, noise weights, figure source tables and immutable input
hashes. The original counters and membership ledger are held in the
controlled experimental archive. {[}Authors to supply an
editor-accessible location and confirm the permanent data-access
statement and repository identifier before submission.{]}

\singlespacing

\hypertarget{code-availability}{%
\subsection{Code availability}\label{code-availability}}

The package contains the independent QR replay, same-data design
comparison, conditional simulations, spectral calculations, covariance
propagation and figure scripts, with software versions and a
reproducible entry point. {[}Authors to confirm the code licence and
permanent repository identifier before submission.{]}

\singlespacing

\hypertarget{author-statements}{%
\subsection{Author statements}\label{author-statements}}

\textbf{Author contributions:} {[}To be completed and approved by all
authors.{]}

\textbf{Acknowledgements and funding:} {[}To be completed by the
authors.{]}

\textbf{Competing interests:} {[}Declaration to be supplied by the
authors.{]}

\textbf{Correspondence:} {[}Corresponding author names and email
addresses to be supplied.{]}

\textbf{AI-assisted preparation:} OpenAI Codex assisted English
drafting, numerical replay and figure preparation. {[}Authors to confirm
their verification of the analyses, citations and final text, and
responsibility for the submitted work.{]}

\clearpage

\singlespacing

\hypertarget{references}{%
\subsection{References}\label{references}}

\begin{enumerate}
\def\labelenumi{\arabic{enumi}.}
\tightlist
\item
  Ludlow, A. D., Boyd, M. M., Ye, J., Peik, E. \& Schmidt, P. O. Optical
  atomic clocks. \emph{Rev.~Mod. Phys.} \textbf{87}, 637--701 (2015).
  \url{https://doi.org/10.1103/revmodphys.87.637}
\item
  Mehlstäubler, T. E., Grosche, G., Lisdat, C., Schmidt, P. O. \&
  Denker, H. Atomic clocks for geodesy. \emph{Rep.~Prog. Phys.}
  \textbf{81}, 064401 (2018).
  \url{https://doi.org/10.1088/1361-6633/aab409}
\item
  Chou, C. W., Hume, D. B., Rosenband, T. \& Wineland, D. J. Optical
  Clocks and Relativity. \emph{Science} \textbf{329}, 1630--1633 (2010).
  \url{https://doi.org/10.1126/science.1192720}
\item
  McGrew, W. F. et al.~Atomic clock performance enabling geodesy below
  the centimetre level. \emph{Nature} \textbf{564}, 87--90 (2018).
  \url{https://doi.org/10.1038/s41586-018-0738-2}
\item
  Grotti, J. et al.~Geodesy and metrology with a transportable optical
  clock. \emph{Nat. Phys.} \textbf{14}, 437--441 (2018).
  \url{https://doi.org/10.1038/s41567-017-0042-3}
\item
  Lisdat, C. et al.~A clock network for geodesy and fundamental science.
  \emph{Nat. Commun.} \textbf{7}, 12443 (2016).
  \url{https://doi.org/10.1038/ncomms12443}
\item
  Riehle, F. Optical clock networks. \emph{Nat. Photon.} \textbf{11},
  25--31 (2017). \url{https://doi.org/10.1038/nphoton.2016.235}
\item
  Lindvall, T. et al.~Coordinated international comparisons between
  optical clocks connected via fiber and satellite links. \emph{Optica}
  \textbf{12}, 843--852 (2025).
  \url{https://doi.org/10.1364/OPTICA.561754}
\item
  Dimarcq, N. et al.~Roadmap towards the redefinition of the second.
  \emph{Metrologia} \textbf{61}, 012001 (2024).
  \url{https://doi.org/10.1088/1681-7575/ad17d2}
\item
  Denker, H. et al.~Geodetic methods to determine the relativistic
  redshift at the level of \(10^{-18}\) in the context of international
  timescales: a review and practical results. \emph{J. Geod.}
  \textbf{92}, 487--516 (2018).
  \url{https://doi.org/10.1007/s00190-017-1075-1}
\item
  Schröder, S., Stellmer, S. \& Kusche, J. Potential and scientific
  requirements of optical clock networks for validating
  satellite-derived time-variable gravity data. \emph{Geophys. J. Int.}
  \textbf{226}, 764--779 (2021).
  \url{https://academic.oup.com/gji/article/226/2/764/6212226}
\item
  Voigt, C., Denker, H. \& Timmen, L. Time-variable gravity potential
  components for optical clock comparisons and the definition of
  international time scales. \emph{Metrologia} \textbf{53}, 1365--1383
  (2016). \url{https://doi.org/10.1088/0026-1394/53/6/1365}
\item
  Hou, L. et al.~Yb/Sr optical clock comparison at \(2.2\times10^{-18}\)
  uncertainty over a 1350-km fiber link. \emph{arXiv}, 2610.34250v1
  (2026). \url{https://arxiv.org/abs/2610.34250v1}
\item
  Jia, Z.-P. et al.~Improved systematic evaluation of a strontium
  optical clock with uncertainty below \(1\times10^{-18}\).
  \emph{Metrologia} \textbf{63}, 025002 (2026).
  \url{https://doi.org/10.1088/1681-7575/ae449e}
\item
  Zhu, Q. et al.~Ytterbium lattice clock with systematic uncertainty of
  \(1.3\times 10^{-18}\) and instability at the \(10^{-19}\) level.
  \emph{arXiv}, 2606.10514v2 (2026).
  \url{https://arxiv.org/abs/2606.10514v2}
\item
  Petit, G. \& Luzum, B. (eds) IERS Conventions (2010). \emph{IERS
  Technical Note} \textbf{36} (2010).
  \url{https://iers-conventions.obspm.fr/content/tn36.pdf}
\item
  CNES. FES2022 (Finite Element Solution) Tidal model (Version 2024).
  CNES \url{https://doi.org/10.24400/527896/A01-2024.004} (2024).
\item
  Farrell, W. E. Deformation of the Earth by surface loads.
  \emph{Rev.~Geophys.} \textbf{10}, 761--797 (1972).
  \url{https://doi.org/10.1029/RG010i003p00761}
\item
  Han, D. \& Wahr, J. The viscoelastic relaxation of a realistically
  stratified earth, and a further analysis of postglacial rebound.
  \emph{Geophys. J. Int.} \textbf{120}, 287--311 (1995).
  \url{https://doi.org/10.1111/j.1365-246x.1995.tb01819.x}
\item
  Stahl, K. et al.~Optical clock network tracks tidal variations of the
  gravity potential. \emph{arXiv}, 2610.39635v1 (2026).
  \url{https://arxiv.org/abs/2610.39635v1}
\item
  Yao, J. et al.~Optical-Clock-Based Time Scale. \emph{Phys.
  Rev.~Applied} \textbf{12}, 044069 (2019).
  \url{https://doi.org/10.1103/PhysRevApplied.12.044069}
\item
  Vincent, A. \& Müller, J. Detection of time variable gravity signals
  using terrestrial clock networks. \emph{Adv. Space Res.} \textbf{73},
  3312--3320 (2024). \url{https://doi.org/10.1016/j.asr.2023.07.058}
\item
  NOAA. Harmonic constituents. \emph{NOAA Tides and Currents}.
  \url{https://tidesandcurrents.noaa.gov/harcon.html?id=8577330}
  (accessed 6 November 2026).
\item
  VanderPlas, J. T. Understanding the Lomb--Scargle periodogram.
  \emph{Astrophys. J. Suppl. Ser.} \textbf{236}, 16 (2018).
  \url{https://doi.org/10.3847/1538-4365/aab766}
\end{enumerate}

\clearpage

\singlespacing\sffamily\fontsize{7}{9}\selectfont

\hypertarget{extended-data}{%
\subsection{Extended Data}\label{extended-data}}

\includegraphics{figures/ed_figure1.pdf}

\textbf{Extended Data Fig. 1 \textbar{} Same-data comparison of
reference-level designs.} \textbf{a,} Weighted mean residual of each
segment after the eight-period joint fit. The additional segment
intercepts make these means zero by construction, so a zero line would
not demonstrate calibration. \textbf{b,} Lunar point estimates after
deleting each observing period under eight-period and 54-segment
designs. Both use the same frozen observations, physical templates,
corrections and noise weights. These diagnostics do not replace
lock/relock or zero-transfer records.

\includegraphics{figures/ed_figure2.pdf}

\textbf{Extended Data Fig. 2 \textbar{} Conditional semidiurnal scans.}
Equal-weight scans of the same 3,299 observations with 54 segment
intercepts. \textbf{a,} M₂ neighbourhood before and after projection of
a sine--cosine pair at S₂. \textbf{b,} S₂ neighbourhood before and after
projection at M₂. Dotted lines show the astronomical target frequencies.
These are descriptive projections on shared observations, not additional
detection probabilities.

\clearpage

\textbf{Extended Data Table 1 \textbar{} Estimator and uncertainty
comparison.} All directly replayed designs use the same 3,299 means,
physical templates and Allan-based weights. Errors are standard errors.
The alternative filtered estimate is separately labelled
author-reported.

\begin{longtable}[]{@{}
  >{\raggedright\arraybackslash}p{(\columnwidth - 8\tabcolsep) * \real{0.2706}}
  >{\raggedright\arraybackslash}p{(\columnwidth - 8\tabcolsep) * \real{0.1765}}
  >{\raggedright\arraybackslash}p{(\columnwidth - 8\tabcolsep) * \real{0.1765}}
  >{\raggedright\arraybackslash}p{(\columnwidth - 8\tabcolsep) * \real{0.1765}}
  >{\raggedright\arraybackslash}p{(\columnwidth - 8\tabcolsep) * \real{0.2000}}@{}}
\toprule\noalign{}
\begin{minipage}[b]{\linewidth}\raggedright
Design
\end{minipage} & \begin{minipage}[b]{\linewidth}\raggedright
Lunar
\end{minipage} & \begin{minipage}[b]{\linewidth}\raggedright
Near\_12h
\end{minipage} & \begin{minipage}[b]{\linewidth}\raggedright
Diurnal
\end{minipage} & \begin{minipage}[b]{\linewidth}\raggedright
Status
\end{minipage} \\
\midrule\noalign{}
\endhead
\bottomrule\noalign{}
\endlastfoot
S54 & 1.000 ± 0.132 & 0.925 ± 0.365 & 0.579 ± 0.538 & Replayed \\
G29 & 0.984 ± 0.124 & 1.093 ± 0.298 & 0.759 ± 0.420 & Replayed \\
P8 & 0.987 ± 0.123 & 1.129 ± 0.288 & 0.678 ± 0.359 & Replayed \\
Filtered, 54 segments & 0.97 ± 0.14 & 0.90 ± 0.29 & 0.35 ± 0.51 &
Author-reported \\
\end{longtable}

\begin{longtable}[]{@{}
  >{\raggedright\arraybackslash}p{(\columnwidth - 10\tabcolsep) * \real{0.1486}}
  >{\raggedleft\arraybackslash}p{(\columnwidth - 10\tabcolsep) * \real{0.1622}}
  >{\raggedleft\arraybackslash}p{(\columnwidth - 10\tabcolsep) * \real{0.1757}}
  >{\raggedleft\arraybackslash}p{(\columnwidth - 10\tabcolsep) * \real{0.1351}}
  >{\raggedleft\arraybackslash}p{(\columnwidth - 10\tabcolsep) * \real{0.1757}}
  >{\raggedleft\arraybackslash}p{(\columnwidth - 10\tabcolsep) * \real{0.2027}}@{}}
\toprule\noalign{}
\begin{minipage}[b]{\linewidth}\raggedright
Group
\end{minipage} & \begin{minipage}[b]{\linewidth}\raggedleft
Estimate
\end{minipage} & \begin{minipage}[b]{\linewidth}\raggedleft
Scaled\_LS
\end{minipage} & \begin{minipage}[b]{\linewidth}\raggedleft
Day\_JK
\end{minipage} & \begin{minipage}[b]{\linewidth}\raggedleft
Period\_JK
\end{minipage} & \begin{minipage}[b]{\linewidth}\raggedleft
Reported\_SE
\end{minipage} \\
\midrule\noalign{}
\endhead
\bottomrule\noalign{}
\endlastfoot
Lunar & 1.0004 & 0.132 & 0.1218 & 0.147 & 0.132 \\
Near-12 h & 0.9251 & 0.2993 & 0.3649 & 0.3531 & 0.3649 \\
Diurnal & 0.5792 & 0.3887 & 0.5381 & 0.7735 & 0.5381 \\
\end{longtable}

\textbf{Extended Data Table 2 \textbar{} Conditional 95\% interval
coverage.} Each scenario has 1,000 simulated realizations per design,
with fixed observed sampling and weights. Values are percentages for the
larger-of-LS-and-day-jackknife rule, using multiplier 1.96. They assess
the specified scenarios, not unknown systematic effects.

\begin{longtable}[]{@{}llrrr@{}}
\toprule\noalign{}
scenario & design & Diurnal & Lunar & Near12h \\
\midrule\noalign{}
\endhead
\bottomrule\noalign{}
\endlastfoot
50\% OU 24 h, period & P8 & 97.9 & 98 & 97.3 \\
50\% OU 24 h, period & S54 & 96.4 & 95.1 & 95.4 \\
50\% OU 6 h, segment & P8 & 94.7 & 95.7 & 96 \\
50\% OU 6 h, segment & S54 & 95.8 & 95.6 & 96.3 \\
Student t (df 5) & P8 & 96.8 & 96.3 & 96.2 \\
Student t (df 5) & S54 & 96.2 & 96.1 & 96.8 \\
Gaussian & P8 & 95.7 & 96.2 & 95.9 \\
Gaussian & S54 & 95.5 & 97.4 & 96.7 \\
\end{longtable}

\textbf{Extended Data Table 3 \textbar{} Astronomical constituent
guide.} Names specify forcing and frequency; one response coefficient
multiplies each complete group with its internal ratios and phases
fixed.

\begin{longtable}[]{@{}lrll@{}}
\toprule\noalign{}
Symbol & Period\_h & Meaning & Group \\
\midrule\noalign{}
\endhead
\bottomrule\noalign{}
\endlastfoot
M2 & 12.421 & Principal lunar semidiurnal & Lunar \\
N2 & 12.658 & Lunar elliptic semidiurnal & Lunar \\
S2 & 12 & Principal solar semidiurnal & Near-12 h \\
K2 & 11.967 & Lunisolar semidiurnal & Near-12 h \\
K1 & 23.934 & Lunisolar diurnal & Diurnal \\
O1 & 25.819 & Principal lunar diurnal & Diurnal \\
P1 & 24.066 & Principal solar diurnal & Diurnal \\
Q1 & 26.868 & Lunar elliptic diurnal & Diurnal \\
\end{longtable}

\end{document}
